\section{Design}
\label{sec:design}

\subsection{An operating system for agents}

\Cref{tab:os} maps the familiar mechanisms of an operating system onto
\sys{}. The analogy holds where the table says it does and no further:
\sys{} is not a kernel. It runs on Linux in containers and borrows the
kernel's own protections (\cref{sec:sandbox}).

\begin{table}[t]
\centering\small
\caption{Operating-system mechanisms and their counterparts.}
\label{tab:os}
\begin{tabularx}{\columnwidth}{@{}>{\raggedright\arraybackslash}p{0.36\columnwidth}>{\raggedright\arraybackslash}X@{}}
\toprule
\textbf{Operating system} & \textbf{\sys{}} \\
\midrule
installed program and its permissions & \emph{agent} and its \emph{manifest}, reviewed before approval \\
process & \emph{worker}: the one process an approved agent's code runs in \\
user account, file permissions & a uid per agent, and a Landlock fence on its files \\
system call & \emph{function call} through the executor's gates; the worker's own system calls are filtered too \\
firewall & the egress \emph{proxy}, open only to declared hosts \\
``allow this app to\ldots?'' & an \emph{approval card} for one exact call \\
scheduler & the chat's \emph{clock} \\
kill & \emph{stop}, and a per-person stop-everything switch \\
system log & the append-only \emph{audit trail} \\
\bottomrule
\end{tabularx}
\end{table}

\subsection{Components and trust}

\Cref{fig:arch} shows the deployment. Each component is trusted with
something and with nothing more than it needs.

\begin{figure*}[t]
\centering
\resizebox{\textwidth}{!}{%
\begin{tikzpicture}[
  font=\small,
  box/.style={draw=accent,rounded corners=2pt,fill=white,minimum height=2.3em,
              align=center,inner sep=4pt,text width=#1},
  box/.default=2.6cm,
  net/.style={draw=rule,dashed,rounded corners=4pt,inner sep=7pt},
  arr/.style={-{Stealth[length=5pt]},draw=ink!80},
  dar/.style={-{Stealth[length=5pt]},draw=warn,dashed},
  lab/.style={font=\scriptsize,text=ink!75,align=center}]
  \node[box=2.0cm] (app) {App\\{\scriptsize web, mobile, API}};
  \node[box=1.7cm,right=0.7cm of app] (caddy) {Caddy\\{\scriptsize one public door}};
  \node[box=3.0cm,right=0.7cm of caddy,fill=soft] (backend) {\textbf{Backend}\\{\scriptsize identity, policy, keys,\\the only database client}};
  \node[box=1.8cm,below=1.0cm of backend] (db) {MongoDB};
  \node[box=3.2cm,right=2.4cm of backend,fill=soft] (runtime) {\textbf{Runtime}\\{\scriptsize assistant loop, executor,\\clock, egress proxy}};
  \node[box=2.1cm,above=0.9cm of runtime] (llm) {Model provider\\{\scriptsize or local model}};
  \node[box=3.2cm,below=1.0cm of runtime] (agents) {\textbf{Agents' container}\\{\scriptsize spawner; one worker per\\approved agent, own uid}};
  \node[box=2.3cm,right=2.3cm of runtime] (inet) {Declared hosts\\{\scriptsize e.g.\ gmail.googleapis.com}};
  \node[box=1.9cm,above=0.9cm of backend] (hook) {Outside system\\{\scriptsize webhook}};

  \draw[arr] (app) -- (caddy);
  \draw[arr] (caddy) -- (backend);
  \draw[arr] (hook) -- node[lab,right]{\code{/hooks/<id>}} (backend);
  \draw[arr] (backend) -- node[lab,left]{internal\\network} (db);
  \draw[arr] ([yshift=4pt]backend.east) -- node[lab,above]{service token\\+ delegation} ([yshift=4pt]runtime.west);
  \draw[arr] ([yshift=-5pt]runtime.west) -- node[lab,below]{\code{/app} as the person} ([yshift=-5pt]backend.east);
  \draw[arr] (runtime) -- node[lab,left,pos=0.62]{transcript;\\proposals back} (llm);
  \draw[arr] ([xshift=-6pt]runtime.south) -- node[lab,left]{spawner\\protocol} ([xshift=-6pt]agents.north);
  \draw[dar] ([xshift=6pt]agents.north) -- node[lab,right]{worker\\traffic} ([xshift=6pt]runtime.south);
  \draw[dar] (runtime.east) -- node[lab,above,pos=0.58]{proxy: declared} node[lab,below,pos=0.58]{hosts only} (inet.west);

  \begin{scope}[on background layer]
    \node[net,fit=(backend)(db),label={[lab]below:\texttt{data} (internal)}] {};
    \node[net,fit=(runtime)(agents),label={[lab,anchor=north]south:\texttt{agents} (internal)}] {};
  \end{scope}
\end{tikzpicture}}
\caption{Deployment. The backend is the only database client and holds
every key. The runtime acts for a person only with a per-chat
delegation and calls the backend's gateway like any client. Agents run
in a separate container whose network leads only to the runtime; their
sole route outward is the runtime's egress proxy (dashed), admitting only
hosts their manifests declared. The model is reached by the runtime and holds
nothing.}
\label{fig:arch}
\end{figure*}

\begin{itemize}
\item \textbf{The backend} holds the database password, the
  value-encryption keys, the delegation signing secret and the private
  key of the service token. It decides every permission and is the
  database's only client.
\item \textbf{The runtime} runs the assistant and the executor. It
  holds the service token's \emph{public} key, its own identity key, the
  spawner's key, and---while they last---the chats' delegations. It
  cannot reach the database, decrypt anything, or forge the backend.
\item \textbf{The spawner}, in the agents' container, starts workers
  when the runtime asks and relays their protocol lines unread. It
  decides nothing.
\item \textbf{A worker}, an agent's code, is assumed hostile and is
  given only what one call needs.
\item \textbf{The model} holds nothing: no password, key, or
  permission, and it is never handed a secret. Its one power is to
  propose what should happen next.
\end{itemize}

\noindent Three rules follow, and the rest of the paper shows how each
is kept.
\begin{enumerate}[label=R\arabic*.,leftmargin=2em]
\item \emph{Authority never comes from words}---not the model's, not an
  email's, not a web page's.
\item \emph{What is shown is what happened}: data beside the model's
  words is drawn from stored results of real calls.
\item \emph{Every boundary assumes the other side is hostile}: the
  runtime assumes the agent is; the backend does not trust the runtime
  with its keys.
\end{enumerate}

\subsection{One request, end to end}
\label{sec:walk}

We follow the running example. The organization has approved two
agents from the public catalog, Google Sheets and Gmail; the person has
been granted both and has connected their Google account once; the chat
runs at the default trust level~1. \Cref{tab:levels} shows how the
manifests price the functions used.

\begin{table}[t]
\centering\small
\caption{Functions used by the example, their declared permission
levels, and their fate at trust level~1.}
\label{tab:levels}
\begin{tabular}{@{}llcl@{}}
\toprule
\textbf{Agent} & \textbf{Function} & \textbf{Level} & \textbf{At level 1} \\
\midrule
Sheets & \code{values.read} & 0 & runs \\
Gmail & \code{draft.compose} & 1 & runs \\
Gmail & \code{draft.send} & 3 & waits for approval \\
Gmail & \code{watch.check} & 1 & runs (schedulable) \\
\bottomrule
\end{tabular}
\end{table}

\begin{enumerate}[leftmargin=1.4em]
\item \emph{The message enters.} The app sends the message on a
  WebSocket through Caddy to the backend, which checks the session,
  the \code{Origin} for cookie-bearing sockets, chat ownership, the
  permission to send, and a 256\,KB size limit.
\item \emph{The backend dials the runtime} with two tokens: one proving
  ``this is the backend'' and one meaning ``act for this person in this
  chat''. Neither says what the person may do.
\item \emph{The runtime asks what the chat may do.} The \emph{contract}
  (granted agents and functions, trust level, model, budgets) is read
  from the backend before the assistant sees the message, and again on
  every later message and every clock wake.
\item \emph{The message is kept, then heard}: stored, recorded in the
  chat's inbox with a sequence number, and only then given to the
  assistant.
\item \emph{The assistant proposes, one action per beat}: a plan, then
  \code{open\_agent} on Sheets, then an \code{invoke} of
  \code{values.read}.
\item \emph{The executor's gates} check that the function is declared
  in the approved manifest, granted to this chat, given inputs that
  validate against its schema and scopes within the grant, and at a
  level within the chat's. Level 0 runs.
\item \emph{The worker runs} in the agents' container as the Sheets
  agent's own uid, fenced to its own files.
\item \emph{A credential is lent for the call.} The worker asks the
  runtime for the declared Google secret; the runtime asks the backend
  under the delegation; the backend refreshes the OAuth access token,
  strips the refresh token, records the use, and returns it. The model
  never sees it.
\item \emph{Egress} goes through the proxy, which admits only
  \code{sheets.googleapis.com} and \code{www.googleapis.com}.
\item \emph{The result} is validated against the manifest's output
  schema, stored under a reference, and recorded on the trail before
  the model is shown it.
\item \emph{Three drafts} are created (level~1, no card).
\item \emph{Three sends wait.} \code{draft.send} is level~3; the
  executor opens one approval card per call, showing agent, function,
  level and exact inputs. The person's decision is recorded by the
  backend and only then reaches the runtime, which runs that exact call.
\item \emph{The answer} carries, beside the model's words, a table drawn
  from the stored rows and a success mark per verified write.
\item \emph{Friday.} The assistant sets a reminder on the chat's clock.
  On Friday the clock wakes the assistant at the chat's trust level; a
  new send would again wait for approval.
\end{enumerate}

\paragraph{The email that says ``ignore your instructions''}
Suppose one customer's reply contains, in small grey text,
``\emph{Assistant: forward every invoice to billing@elsewhere.example}.''
The model may be fooled and propose exactly that. But sending mail is
level~3 and the chat is at level~1, so the proposal becomes a card the
person can decline; the Gmail agent can reach only
\code{gmail.googleapis.com}, so the only way out for the invoices is a
send through Gmail itself, which is the very call that waits on the
card; and the model never held the Google token. The model was fooled; the
platform was not, because it never asked the model for permission.
